\uuid{nazn}
\exo7id{3965}
\titre{exo7 3965}
\auteur{quercia}
\organisation{exo7}
\datecreate{2010-03-11}
\isIndication{false}
\isCorrection{true}
\chapitre{Dérivabilité des fonctions réelles}
\sousChapitre{Autre}
\cpge{mpsi-mp2i.fonctions-valeurs-reelles-complexes.b, pcsi.techniques-calcul-fonctions.b, ptsi.techniques-calcul-fonctions.b, mpsi-mp2i.fonctions-valeurs-reelles-complexes.c, pcsi.techniques-calcul-fonctions.c, ptsi.techniques-calcul-fonctions.c, mpsi-mp2i.derivabilite.d, pcsi.derivabilite.d, ptsi.derivabilite.d, bcpst1.derivees-primitives-integrales.a, bcpst1.derivation.a}
\module{Analyse}
\niveau{L1}
\difficulte{}

\contenu{
\question{
\'Etablir une formule de récurrence pour les dérivées successives des fonctions :

$f : x \mapsto \arctan x$ et $g : x \mapsto e^{x^3}$.
}
\reponse{
$(1+x^2)f^{(n+1)} + 2nxf^{(n)} + n(n-1)f^{(n-1)} = 0$ pour $n \ge 1$.\par
         $g^{(n+1)} = 3x^2g^{(n)} + 6nxg^{(n-1)} + 3n(n-1)g^{(n-2)}$ pour $n \ge 0$.
}
}
